\capitulosinnumero{Glosario}

\section*{Siglas y Acrónimos}
\begin{glosario}
  \termino{API}{Interfaz de Programación de Aplicaciones (\textit{Application Programming Interface}).}
  \termino{AS-IS / TO-BE}{Proceso actual y proceso propuesto.}
  \termino{AUC}{Área bajo la curva (\textit{Area Under the Curve}).}
  \termino{CORS}{Intercambio de recursos entre orígenes (\textit{Cross-Origin Resource Sharing}).}
  \termino{DHS}{Encuestas Demográficas y de Salud (\textit{Demographic and Health Surveys}).}
  \termino{ENDI}{Encuesta Nacional sobre Desnutrición Infantil.}
  \termino{ENSANUT}{Encuesta Nacional de Salud y Nutrición.}
  \termino{FN, FP, VN, VP}{Falsos negativos, falsos positivos, verdaderos negativos y verdaderos positivos.}
  \termino{HTTP}{Protocolo de transferencia de hipertexto (\textit{HyperText Transfer Protocol}).}
  \termino{IC}{Intervalo de confianza.}
  \termino{INEC}{Instituto Nacional de Estadística y Censos.}
  \termino{JSON}{Notación de objetos de JavaScript (\textit{JavaScript Object Notation}).}
  \termino{ML}{Aprendizaje automático (\textit{Machine Learning}).}
  \termino{ODS}{Objetivos de Desarrollo Sostenible.}
  \termino{OMS}{Organización Mundial de la Salud.}
  \termino{OSEMN}{\textit{Obtain, Scrub, Explore, Model, iNterpret}.}
  \termino{PR-AUC}{Área bajo la curva de precisión--sensibilidad (\textit{average precision}).}
  \termino{REST}{Transferencia de Estado Representacional (\textit{Representational State Transfer}).}
  \termino{ROC}{Característica operativa del receptor (\textit{Receiver Operating Characteristic}).}
  \termino{SHAP}{Explicaciones aditivas de Shapley (\textit{SHapley Additive exPlanations}).}
\end{glosario}

\section*{Glosario de Términos}
\begin{glosario}
  \termino{Bootstrap}{Técnica que remuestrea con reemplazo los datos de prueba para estimar la variabilidad de una métrica y su intervalo de confianza.}
  \termino{Brier (error de)}{Promedio del cuadrado de la diferencia entre la probabilidad predicha y el resultado observado; mide la calibración y la precisión de las probabilidades.}
  \termino{Calibración}{Grado en que las probabilidades predichas coinciden con las frecuencias observadas.}
  \termino{Clasificación}{Tarea de asignar cada registro a una categoría; aquí, con o sin desnutrición crónica.}
  \termino{Desnutrición crónica}{Retraso del crecimiento en talla para la edad, asociado a deficiencias nutricionales prolongadas.}
  \termino{Desbalance de clases}{Situación en que una categoría tiene muchos menos registros que otra.}
  \termino{Diseño transversal}{Estudio en el que cada unidad se mide una sola vez, por lo que no permite anticipar eventos futuros.}
  \termino{Especificidad}{Proporción de casos negativos que el modelo identifica correctamente.}
  \termino{Estrato}{Variable de diseño muestral que agrupa unidades de la encuesta; no describe al niño.}
  \termino{Exactitud balanceada}{Promedio de la sensibilidad y la especificidad.}
  \termino{Factor de expansión}{Peso que indica a cuántas unidades de la población representa un registro de la encuesta.}
  \termino{Fuga de información}{Uso, durante el entrenamiento, de datos que revelan la variable objetivo o la partición de prueba.}
  \termino{One-hot}{Codificación de una variable categórica en columnas binarias, una por categoría.}
  \termino{Pipeline}{Secuencia de pasos de preprocesamiento y modelado que se ajusta como una sola unidad.}
  \termino{Precisión}{Proporción de positivos predichos que son casos reales.}
  \termino{Prevalencia}{Proporción de casos en una población o muestra en un momento dado.}
  \termino{Random Forest}{Modelo que promedia muchos árboles de decisión entrenados con muestras aleatorias.}
  \termino{Sensibilidad}{Proporción de casos positivos que el modelo identifica correctamente.}
  \termino{Sobreajuste}{Situación en que un modelo se ajusta al entrenamiento mucho mejor que a datos nuevos.}
  \termino{Tamizaje}{Aplicación de una prueba a una población para detectar posibles casos.}
  \termino{Umbral de decisión}{Probabilidad a partir de la cual el modelo asigna la clase positiva.}
  \termino{Validación cruzada}{Procedimiento que divide el entrenamiento en particiones para ajustar y evaluar sin usar la prueba.}
  \termino{Validación externa}{Evaluación con datos de otra fuente, población o periodo distintos de los usados para entrenar.}
\end{glosario}
